% L'architettura di un progetto di intelligenza artificiale in Halue: l'applicazione dell'agente
% sta tra l'utente, il modello e i servizi esterni da cui arrivano dati e azioni.
% In bordeaux le parti sviluppate per ogni progetto.
% Il disegno è largo quanto il testo (150 mm): si include con \input, senza ridimensionarlo.

\begin{tikzpicture}[
	every node/.style={text=FiguraTesto},
	blocco/.style={draw=FiguraTenue, fill=white, rounded corners=3pt, text width=#1-2mm,
		minimum width=#1, minimum height=11mm, align=center, inner sep=1.5pt,
		font=\fontsize{7.5}{9}\selectfont,
		execute at begin node={\hyphenpenalty=10000\exhyphenpenalty=10000}},
	blocco/.default=40mm,
	interno/.style={draw=FiguraAccento, line width=0.8pt},
	doppia/.style={freccia, {Stealth[length=1.8mm,width=1.4mm]}-{Stealth[length=1.8mm,width=1.4mm]}},
	etichetta/.style={font=\fontsize{7.5}{9}\selectfont, text=FiguraTenue, inner sep=1pt,
		align=center},
]
	% L'applicazione dell'agente e i servizi esterni, come contenitori sullo sfondo.
	\begin{scope}[on background layer]
		\path[partizione] (32mm, -24mm) rectangle (98mm, -76mm);
		\path[intestazione] (32mm, -24mm) rectangle (98mm, -30mm);
		\node[titolo] at (65mm, -27mm) {Applicazione dell’agente};
		\path[partizione] (106mm, -24mm) rectangle (150mm, -76mm);
		\path[intestazione] (106mm, -24mm) rectangle (150mm, -30mm);
		\node[titolo] at (128mm, -27mm) {Servizi esterni};
	\end{scope}

	% Il modello, fuori dall'applicazione che lo usa.
	\node[blocco=66mm] (llm) at (65mm, -6.5mm)
		{\voce{Modello linguistico (LLM)}{servizio del fornitore, di solito Claude}};

	% Dentro l'applicazione.
	\node[blocco=54mm, interno] (istruzioni) at (65mm, -38.5mm)
		{\voce{Istruzioni}{compito, regole, tono del marchio}};
	\node[blocco=54mm, interno] (strumenti) at (65mm, -52.5mm)
		{\voce{Strumenti interni}{ricerche e azioni su misura del progetto}};
	\node[blocco=54mm] (client) at (65mm, -66.5mm)
		{\voce{\emph{Client} MCP}{scopre e invoca gli strumenti dei \emph{server}}};

	% Fuori dall'applicazione.
	\node[blocco=22mm, minimum height=16mm] (utente) at (11mm, -52.5mm)
		{\voce{Utente}{scrive in linguaggio naturale}};
	\node[blocco=40mm] (ucp) at (128mm, -38.5mm)
		{\voce{Negozio con UCP}{carrello, \emph{checkout}, pagamento}};
	\node[blocco=40mm] (api) at (128mm, -52.5mm)
		{\voce{API della piattaforma}{Shopify, Salesforce, gestionali}};
	\node[blocco=40mm] (mcp) at (128mm, -66.5mm)
		{\voce{\emph{Server} MCP}{della piattaforma o di terze parti}};

	% Il ciclo tra applicazione e modello.
	\draw[freccia] (58mm, -24mm) -- (58mm, -12mm);
	\draw[freccia] (72mm, -12mm) -- (72mm, -24mm);
	\node[etichetta, anchor=east] at (56.5mm, -18mm) {conversazione, istruzioni,\\descrizione degli strumenti};
	\node[etichetta, anchor=west] at (73.5mm, -18mm) {testo, oppure\\strumento da invocare};

	% Collegamenti con l'esterno.
	\draw[doppia] (utente.east) -- (32mm, -52.5mm);
	\draw[freccia] (strumenti.east) -- (api.west);
	\draw[freccia] (102mm, -52.5mm) |- (ucp.west);
	\fill[FiguraTenue] (102mm, -52.5mm) circle (0.5mm);
	\draw[freccia] (client.east) -- (mcp.west);
\end{tikzpicture}
